\section{Methodology: What We Got Wrong and How We Fixed It}
\label{sec:methodology}

This section is not a methods description; it is an audit. Five defects in our
own comparison harness biased the measurements before any architecture
conclusion was drawn. Four of the five biased against the autoregressive (AR)
baseline. We found all five not by running new experiments but by re-reading
committed artifacts, and the fixes moved headline numbers by more than the
effects we were reporting. We therefore report the defects as a first-class
result rather than as a footnote to the experimental setup.

\subsection{Pinned stack and provenance}

Every measurement from issue~\#57 onward uses a single pinned stack
(Table~\ref{tab:pinned-stack}): vLLM~\vllmVersion{} on an \gpuModel{}. A
shared provenance helper, which no runner can bypass, writes the model id,
revision, runtime versions, GPU, dtype, and every decoding knob into each
result JSON from the object that actually loaded the model. A build check
(\texttt{paper/verify\_numbers.py}) fails when a cited artifact lacks the
block. The helper exists because three result files recorded the wrong model
name for the diffusion arm (defect~5, below). On this stack, run-to-run
coefficient of variation at $L=\canvasLength$ is AR~0.05\% and DG~0.25\%
(five repeats), so the effects reported below are not noise.

\begin{table}[t]
\centering
\caption{Pinned stack. All numbers in this paper are measured here unless
explicitly labelled otherwise.}
\label{tab:pinned-stack}
\begin{tabular}{ll}
\toprule
Component & Version \\
\midrule
vLLM & \vllmVersion \\
torch & \torchVersion \\
GPU & \gpuModel \\
\bottomrule
\end{tabular}
\end{table}

\subsection{Nine-field scoring}

We score extraction quality on 9 fields, not 12. Three ``semantic'' fields
have degenerate gold labels that score $\sim\!1.0$ for an empty or wrong
prediction, adding a bias that inflates whichever arm happens to interact
badly with them. Dropping them moves the comparison from AR-favoured
(12-field: AR~\arThreeBFtwelve{} vs DG~\dgTwelveF{}) to comparable (9-field:
AR~\arThreeBFnine{} vs DG~\dgNineF{}, spread~$<\!0.01$). The 12-field numbers
are retained for provenance but are not cited in comparison tables.

\subsection{The five defects}

Table~\ref{tab:defects} lists the defects, the magnitude each caused, and the
direction of the bias. All five were found by re-reading committed artifacts;
none required a new GPU run.

\begin{table}[t]
\centering
\caption{Five harness defects. Four of five biased against the AR baseline.
``Before'' is the number as originally recorded; ``After'' is the corrected
value on the pinned stack.}
\label{tab:defects}
\begin{tabular}{clll}
\toprule
\# & Defect & Magnitude (before $\to$ after) & Bias \\
\midrule
1 & Decoding config       & F1 \defectOneFBefore{} $\to$ \defectOneFAfter{}                        & AR \\
2 & Serving parity        & \defectTwoBefore{} $\to$ \defectTwoAfter{} ms                         & AR \\
3 & Parser policy         & F1 \defectThreeBefore{} $\to$ \defectThreeAfter{}                     & AR \\
4 & Metric defaults       & \defectFourBefore\% $\to$ \defectFourAfter\% completeness            & AR \\
5 & Provenance            & 3 files mislabelled                                                   & neutral \\
\bottomrule
\end{tabular}
\end{table}

\paragraph{Decoding config (defect~1).}
HuggingFace \texttt{generate()} honours \texttt{repetition\_penalty} even under
\texttt{do\_sample=False}. Qwen2.5's checkpoint default (1.1 at 1.5B, 1.05 at
3B/7B) emptied the \texttt{host} field on 80/100 events. DiffusionGemma's
hand-written denoising loop has no such default. Setting every knob explicitly
on both arms moved AR-1.5B F1 from \defectOneFBefore{} to \defectOneFAfter{}.

\paragraph{Serving parity (defect~2).}
The AR arm ran a naive HuggingFace generate loop while the diffusion arm ran a
production-grade engine. Putting AR on vLLM moved its p50 from
\defectTwoBefore{} to \defectTwoAfter{}~ms ($4.5\times$) and reversed the
direction of the headline speed claim.

\paragraph{Parser policy (defect~3).}
Forced to over-generate at $L=512$, 96\% of AR outputs contained more than one
JSON object. The old parser took the span from the first brace to the
\emph{last} brace, scoring the most-drifted copy of each field. Re-scoring
identical saved generations with first-object-wins moved AR-3B $L=512$ F1 from
\defectThreeBefore{} to \defectThreeAfter{} ($+0.23$). DiffusionGemma rambles
1--3\% and is parser-invariant.

\paragraph{Metric defaults (defect~4).}
The completeness metric compared \texttt{rec.get(f)} against the defaults
\texttt{""}/\texttt{0.0}/\texttt{False}, which are structurally valid values,
not missing keys. AR stops early and emits falsy defaults; DiffusionGemma
always fills all fields. Distinguishing absent key from present-but-falsy moved
AR-3B completeness from \defectFourBefore\% to \defectFourAfter\% and refuted
the finding it had supported.

\paragraph{Provenance (defect~5).}
A CLI default survived into three result files, recording an AR model name for
diffusion runs. Neutral on any comparison, but it made audit impossible until
correction sidecars (\texttt{*.provenance.json}) were written.

\subsection{The asymmetry}

The four AR-biased defects share a pattern: each lives in the
\emph{infrastructure the AR model runs on}---its checkpoint defaults, its
serving path, its over-generation behaviour, its early-stopping behaviour---not
in the model itself. A harness written by someone hoping the new architecture
wins fails asymmetrically, because the baseline inherits the defaults, the
naive serving path, and the production-grade infrastructure of the system it is
trying to beat. The checklist (\texttt{CHECKLIST.md}) generalises the five
checks so they can be applied before publication; each takes under five
minutes and catches a failure that moved a headline number by more than the
effect it was reporting.
